% Owner: root. classification_e and boundary_d full-source referee PASS.
\begin{proposition}[E.2.17 (Proved): persistence of the normal-crossing cusp]
In the smooth fixed-pair specialization of E.2.10,
the cusp divisor is a relative normal-crossing divisor
over the parameter disc.  More precisely, after
shrinking that disc, at every point of the total cusp
divisor there are holomorphic coordinates
$(z_1,z_2,z_3,s)$, with $s$ the original parameter,
in which the cusp coordinate of the intrinsic map is
one of
\[
 q=z_1,\qquad q=z_1z_2,\qquad q=z_1z_2z_3.
 \tag{E.2.17.1}
\]
There are exactly two relative triple-point sections.
Both lie on the fixed rational curve of E.2.14.
In particular the normalization of the total cusp
divisor is smooth and proper over the parameter
disc, and its central fibre is the normalization
of the central cusp.
\end{proposition}
\begin{proof}
By E.2.14--E.2.16 the three labelled critical
values remain distinct, the cusp has multiplicity
one, and its reduced central fibre is irreducible.
We normalize its target section to $q=0$.
On a sufficiently small target disc, zero is the
only critical value on every parameter fibre.
The extended field $\delta$ is zero precisely on
the relative fixed rational curve $\mathcal C$,
and the two normal characters there are $+1,-1$.
For $s\ne0$ the two triple points of the prescribed
cusp belong to $\mathcal C_s$; outside that curve
the cusp is smooth or has an ordinary double crossing.

We first record exactly the local singularity result
we use.  If $h_s(a,b)$ is a holomorphic family and
$h_0=0$ is a reduced plane-curve germ with an
isolated critical point, choose a sufficiently
small representative of that germ.  The sum of
the multiplicities of all nearby critical points
of $h_s$, regardless of their critical values,
is $\mu(h_0)$.  One can see this directly from
the analytic critical scheme
\[
 \mathbb C\{a,b,s\}/
       ((h_s)_a,(h_s)_b).
\]
After taking a sufficiently small representative,
it is finite over the parameter disc.  Its two
equations are a regular sequence, and its
central fibre is zero-dimensional.  It is thus
Cohen--Macaulay of dimension one, with $s$ a
nonzerodivisor, and is finite flat over the
disc.  Its degree is the central Jacobian-algebra
dimension, namely $\mu(h_0)$.

The non-splitting theorem of
\citet[Theorem 2 and its corollary, pp.~274--275]{Lazzeri1973Root}
says that if a deformation of an isolated
hypersurface singularity has more than one
singular point in a single deformed fibre,
their summed Milnor numbers are strictly
smaller than the original Milnor number.
The strict numerical inequality is the
conclusion of the cited corollary's proof.
Consequently, if all the critical points
just counted have critical value zero,
there can be only one such point.  This
application uses a flat deformation of
plane hypersurfaces: the map
$(a,b,s)\mapsto(h_s(a,b),s)$ is flat near
the central germ because its two base
parameters form a regular sequence in
the smooth source.  No assertion of
separatedness of compact threefolds enters
this local argument.

Let $p$ be a central cusp point outside
$\mathcal C_0$.  A relative holomorphic
flow box for the nonzero field $\delta$
gives coordinates $(a,b,v,s)$ in which
$\delta=\partial_v$.  Invariance of the
intrinsic map gives
\[
 q=h_s(a,b).
\]
The central equation is reduced, by
E.2.16.  A reduced plane-curve germ is
either smooth or has an isolated critical
point.  In the smooth case the implicit
function theorem already gives the first
form in (E.2.17.1).

In the singular case every critical point
of $h_s$ in a small representative lifts
to a critical point of the actual map
on $X_s$.  All its critical values are
therefore zero.  For $s\ne0$ each such
point is a node, since the flow box is
disjoint from the fixed curve and hence
from the two triple points.  Conservation
and non-splitting give one node and
$\mu(h_0)=1$.  Thus the central Hessian
is nondegenerate.  The implicit function
theorem gives its unique critical-point
section, whose critical value is zero.
The holomorphic Morse lemma with
parameters then gives $q=ab$ throughout
this neighbourhood, with $s$ unchanged.

It remains to treat the fixed curve.
Holomorphic coordinates obtained by
averaging over the extended compact
circle may be chosen in the form
$(z,x,y,s)$, with $\mathcal C$ defined
by $x=y=0$ and with circle weights
$0,+1,-1,0$.  The original parameter
$s$ is retained by this averaging.
The holomorphic invariant function $q$
vanishes on $\mathcal C$, so its
convergent power series has the form
\[
 q=F_s(T,z)=T g_s(T,z),\qquad T=xy.
 \tag{E.2.17.2}
\]
The coefficient
$Q_s(z)=g_s(0,z)$ is intrinsically a
section of the dual of the product of
the two normal weight lines along
$\mathcal C_s$.  By E.2.14 this is
$\mathcal O_{\mathbb P^1}(2)$.

The central section $Q_0$ is not
identically zero.  Otherwise the
central equation in every such chart
would be divisible by $x^2y^2$,
contradicting reducedness of the
central cusp.  Its zero divisor
therefore has degree two.  Suppose
that both zeros coincide at a point
$z=0$.  The plane equation
$F_0(T,z)=Tg_0(T,z)$ is reduced:
a repeated factor of $g_0$ would
remain repeated after $T=xy$, and
$T$ does not divide $g_0$ because
$Q_0$ is not identically zero.
Thus $F_0$ has an isolated plane
critical point there.

For $s\ne0$ the two zeros of $Q_s$
are the two distinct triple points
of the prescribed cusp.  Under the
supposed collision they both lie
in the chosen small plane
representative.  They are distinct
critical points of $F_s$ with value
zero.  Every other critical point
of $F_s$ also has value zero:
this is automatic at $T=0$, and
at $T\ne0$ it follows by lifting
to $xy=T$, where the quotient
map $(x,y,z)\mapsto(T,z)$ is a
submersion.  All sufficiently small
$T$ have such a lift in the original
coordinate chart.  A nonzero
critical value would contradict
the actual critical-value set of
the intrinsic map.

Conservation of all critical
multiplicities and Lazzeri's
non-splitting theorem now give a
contradiction to the two distinct
critical points.  Hence $Q_0$ has
two simple zeros.  They extend
as two disjoint holomorphic
sections over the parameter disc.
Where $Q_0\ne0$, the factor
$g_s(xy,z)$ is a unit and can
be absorbed into $y$, giving
$q=xy$.  At a simple zero,
$\partial g_0/\partial z\ne0$;
replace $z$ by $g_s(xy,z)$ to
obtain $q=xyz$, again keeping
$s$ unchanged.  This proves
all local forms in (E.2.17.1).

Properness makes the central cusp
compact, so finitely many such
neighbourhoods cover it and a
single shrinking of the parameter
disc suffices.  The normalization
of a relative normal-crossing
divisor separates its local
branches.  In the coordinates
above every normalized branch
is smooth over $s$, and these
descriptions agree on overlaps.
Normalization is finite for
complex analytic spaces; its
composition with the proper
cusp map is therefore proper.
Its fibres are precisely the
normalizations of the cusp
fibres in these local charts.
This proves the final assertions.
\end{proof}

\begin{remark}[Literature and remaining identification]
Non-splitting is a known theorem of
Lazzeri, not a new singularity-theoretic
claim.  We opened the original paper
and use its strict numerical
corollary.  The application combines
that theorem with the actual extended
circle action, reducedness, and the
absence of extra critical values
proved above.  Searches for
non-splitting, Milnor-number
conservation and degenerations of
torus-fibred complex six-spheres
did not locate this application;
no priority is claimed.  The
proposition supplies relative
normal-crossing charts and a
simultaneous normalization.  It
does not by itself identify the entire
analytic cusp neighbourhood.  The
prescribed marked Mumford completion
is recovered by E.2.18--E.2.20.
\end{remark}
